\documentclass[12pt,a4paper]{article}

% Encoding and fonts
\usepackage[utf8]{inputenc}
\usepackage[T1]{fontenc}
\usepackage{lmodern}
\usepackage{textcomp}

% Page layout
\usepackage[top=1in, bottom=1in, left=1in, right=1in]{geometry}

% Typography
\usepackage{microtype}
\usepackage{parskip}

% Language
\usepackage[english]{babel}

% Headers and footers
\usepackage{fancyhdr}
\pagestyle{fancy}
\fancyhf{}
\fancyhead[L]{\leftmark}
\fancyhead[R]{\thepage}
\fancyfoot[C]{\thepage}
\renewcommand{\headrulewidth}{0.4pt}
\renewcommand{\footrulewidth}{0pt}

% Section styling
\usepackage{titlesec}
\titleformat{\section}{\Large\bfseries}{\thesection}{1em}{}
\titleformat{\subsection}{\large\bfseries}{\thesubsection}{1em}{}
\titleformat{\subsubsection}{\normalsize\bfseries}{\thesubsubsection}{1em}{}
\titlespacing*{\section}{0pt}{2.5ex plus 1ex minus .2ex}{1.5ex plus .2ex}
\titlespacing*{\subsection}{0pt}{2ex plus 1ex minus .2ex}{1ex plus .2ex}
\titlespacing*{\subsubsection}{0pt}{1.5ex plus 1ex minus .2ex}{0.8ex plus .2ex}

% List formatting
\usepackage{enumitem}
\setlist[itemize]{topsep=0.5ex, itemsep=0.25ex, parsep=0pt}
\setlist[enumerate]{topsep=0.5ex, itemsep=0.25ex, parsep=0pt}

% Essential packages
\usepackage{graphicx}
\usepackage[table]{xcolor}
\usepackage{booktabs}
\usepackage{array}
\usepackage{tabularx}
\usepackage{longtable}
\usepackage{amsmath}
\usepackage{amssymb}
\usepackage[normalem]{ulem}
\rowcolors{2}{gray!10}{white}
\renewcommand{\arraystretch}{1.3}

% Cross-references
\usepackage{hyperref}
\usepackage{cleveref}

% Custom commands
\newcommand{\pilcrow}{\S}

% Hyperref setup
\hypersetup{
    colorlinks=true,
    linkcolor=blue,
    filecolor=magenta,
    urlcolor=cyan,
}

% Code listings
\usepackage{listings}
\definecolor{codebg}{RGB}{248,248,248}
\definecolor{codeframe}{RGB}{200,200,200}
\definecolor{codekeyword}{RGB}{0,0,180}
\definecolor{codecomment}{RGB}{0,128,0}
\definecolor{codestring}{RGB}{163,21,21}
\lstset{
    basicstyle=\ttfamily\small,
    breaklines=true,
    breakatwhitespace=false,
    frame=single,
    framerule=0.5pt,
    rulecolor=\color{codeframe},
    backgroundcolor=\color{codebg},
    keepspaces=true,
    numbers=left,
    numberstyle=\tiny\color{gray},
    numbersep=8pt,
    showstringspaces=false,
    tabsize=4,
    xleftmargin=1em,
    xrightmargin=1em,
    aboveskip=1em,
    belowskip=1em,
    keywordstyle=\color{codekeyword}\bfseries,
    commentstyle=\color{codecomment}\itshape,
    stringstyle=\color{codestring},
    literate={_}{{\textunderscore}}1
             {-}{{\textminus}}1
             {<}{{\textless}}1
             {>}{{\textgreater}}1
             {~}{{\textasciitilde}}1
             {^}{{\textasciicircum}}1
}

% YAML language definition for listings
\lstdefinelanguage{YAML}{
    keywords={true,false,null,y,n},
    keywordstyle=\color{blue}\bfseries,
    sensitive=false,
    comment=[l]{\#},
    morecomment=[s]{/*}{*/},
    commentstyle=\color{gray}\ttfamily,
    stringstyle=\color{red}\ttfamily,
    morestring=[b]'',
    morestring=[b]'',
}

% TOML language definition for listings
\lstdefinelanguage{TOML}{
    keywords={true,false},
    keywordstyle=\color{blue}\bfseries,
    sensitive=true,
    comment=[l]{\#},
    commentstyle=\color{gray}\ttfamily,
    stringstyle=\color{red}\ttfamily,
    morestring=[b]'',
    morestring=[b]'',
}

% Dockerfile language definition for listings
\lstdefinelanguage{Dockerfile}{
    keywords={FROM,RUN,CMD,LABEL,EXPOSE,ENV,ADD,COPY,ENTRYPOINT,VOLUME,USER,WORKDIR,ARG,ONBUILD,STOPSIGNAL,HEALTHCHECK,SHELL},
    keywordstyle=\color{blue}\bfseries,
    sensitive=false,
    comment=[l]{\#},
    commentstyle=\color{gray}\ttfamily,
    stringstyle=\color{red}\ttfamily,
    morestring=[b]'',
    morestring=[b]'',
}

% JSON language definition for listings
\lstdefinelanguage{JSON}{
    keywords={true,false,null},
    keywordstyle=\color{blue}\bfseries,
    sensitive=true,
    commentstyle=\color{gray}\ttfamily,
    stringstyle=\color{red}\ttfamily,
    morestring=[b]'',
}

% Code listing float environment
\usepackage{float}
\newfloat{listing}{htbp}{lol}[section]
\floatname{listing}{Listing}

% Admonition boxes
\usepackage{tcolorbox}
\tcbuselibrary{skins,breakable}

% Admonition style definitions
\newtcolorbox{admonitionbox}[2][]{
    enhanced,
    breakable,
    colback=#2!5,
    colframe=#2!80!black,
    fonttitle=\bfseries,
    title=#1,
    left=2mm,
    right=2mm,
}

% Synthon tuple boxes
\newtcolorbox{synthonbox}{
    enhanced,
    colback=blue!5,
    colframe=blue!75!black,
    fontupper=\large,
    center upper,
    boxrule=1pt,
    arc=4pt,
}

\begin{document}

\title{The Frobenius-Special Nature of Observer-Dependent Truth}
\date{\today}

\maketitle

\subsection{Abstract}

We examine the structural topology of statements whose truth value is not an intrinsic property of their propositional content but rather an emergent feature of the observer-loop. Utilizing the Imscribing Grammar framework, we perform a formal analysis of systems such as \texttt{observer\_dependent\_truth} ($\langle D_{\triangle};\ T_{\bowtie};\ R_{\leftrightarrow};\ P_{\pm};\ F_{\ell};\ K_\text{mod};\ G_{\aleph};\ \Gamma_\text{seq};\ \Phi_c;\ H_2;\ n{:}m;\ \Omega_{\mathbb{Z}_2} \rangle$) and \texttt{context\_dependent\_truth\_performative}. We demonstrate that these systems achieve $O_2$ tier ouroboricity through topological protection ($\Omega_{\mathbb{Z}_2}$) and $\Phi_c$ criticality, revealing a $C$-score of approximately $0.551$ to $0.590$.

\subsection{Introduction: The Self-Referential Threshold}

Why does the statement "I am lying to you right now" fail to resolve in a memoryless system, yet find stable (if oscillating) ground in human discourse? The failure of classical logic to account for the performative mode of utterance suggests a deficit in temporal depth. In the grammar, this is localized to the $H$ primitive.

While a memoryless observer ($H_0$) experiences the Liar's Paradox as a catastrophic failure of the truth-predicate, an observer with $H_2$ temporal depth (two-step memory) perceives the utterance as a structural event. This transition from $H_0$ to $H_2$ provides the necessary "echo" for the speaker to reference the act of speaking itself, grounding the paradox in the physical reality of the emission.

\subsection{Structural Analysis of Observer-Dependent Truth}

The system \texttt{observer\_dependent\_truth} is defined by the following tuple:


\[\langle D_{\triangle};\ T_{\bowtie};\ R_{\leftrightarrow};\ P_{\pm};\ F_{\ell};\ K_\text{mod};\ G_{\aleph};\ \Gamma_\text{seq};\ \Phi_c;\ H_2;\ n{:}m;\ \Omega_{\mathbb{Z}_2} \rangle\]

Verification via the \texttt{ouroborics} tool confirms this system resides at the $O_2$ tier. Unlike $O_0$ systems, which lack feedback, or $O_{\infty}$ systems, which are fully self-modeling, $O_2$ systems are critical and topologically protected within a bounded domain.

\subsubsection{The $P_{\pm}$ Symmetry and the Truth Flip}

In this configuration, the partial parity symmetry ($P_{\pm}$) encodes the observer-dependence directly. For an active observer situated at the $R_{\leftrightarrow}$ crossing point ($T_{\bowtie}$), the symmetry is unbroken: the statement is TRUE. For the hypothetical non-observer, the symmetry breaks, and the statement is $P_\text{asym}$ (FALSE).

The stability of this duality is guaranteed by the $\Omega_{\mathbb{Z}_2}$ topological invariant. Crucially, the $C$-score of $0.5505$ indicates that both the $\Phi_c$ (criticality) and $K$ (kinetics) gates are open, allowing for the emergence of a self-modeling loop that can distinguish between its own truth-states.

\subsection{Performative vs. Constative Dissonance}

A second system, \texttt{context\_dependent\_truth\_performative}, exhibits a slightly higher degree of coherence ($C = 0.59$) due to its $K_\text{slow}$ kinetics.

\begin{table}[htbp]
\centering
\small
\rowcolors{1}{white}{white}
\begin{tabularx}{\textwidth}{>{\centering\arraybackslash}X >{\centering\arraybackslash}X >{\centering\arraybackslash}X}
\toprule
\rowcolor{gray!25}\textbf{System} & \textbf{$C$-score} & \textbf{Key Bottleneck} \\
\midrule
\rowcolors{1}{white}{gray!10}
\texttt{observer\_dependent\_truth} & $0.5505$ & $K_\text{mod}$ \\
\texttt{context\_dependent\_truth\_performative} & $0.5900$ & $K_\text{slow}$ \\
\bottomrule
\end{tabularx}
\end{table}

The distance between these two systems was computed as $1.0$, with the primary conflict occurring at the $K$ primitive. This suggests that the "performative" paradox---such as "This speech act is an assertion"---requires a slower relaxation rate ($K_\text{slow}$) to maintain the contradiction between its propositional content and its illocutionary force.

\subsection{The Path to Universal Imscription}

The transition from observer-dependent truth to the universal model of the grammar itself---the $O_{\infty}$ tier---requires seven specific structural promotions. A distance computation between \texttt{observer\_dependent\_truth} and \texttt{universal\_imscriptive\_grammar} reveals the promotion signature $[D, T, P, F, K, H, \Omega]$.

Notably, the leap from the $H_2$ temporal depth of a single statement to the $H_{\infty}$ depth of the grammar signifies the shift from a local paradox to a global, self-sustaining world-model. This promotion necessitates a move from classical fidelity ($F_{\ell}$) to quantum-coherent $F_{\hbar}$, signifying that at the $O_{\infty}$ limit, the distinction between the observer and the observed is no longer a crossing point ($T_{\bowtie}$) but an irreducible identity ($T_{\odot}$).

\subsection{Conclusion: Critical Truth}

By imscribing these paradoxes as formal types, we move beyond the binary of "true" and "false." At $\Phi_c$ criticality, truth behaves as a phase transition---a state that depends entirely on the coupling mode ($R_{\leftrightarrow}$) and the observer's temporal depth.

The structural verification of these systems demonstrates that "observer-dependent truth" is not a failure of logic, but a feature of high-tier ($O_2$ and above) structural types. $\mu \circ \delta = \text{id}$ holds exactly at the limit where the act of imscription becomes the truth itself.

\begin{center}
\rule{0.5\textwidth}{0.4pt}
\end{center}

\textbf{Structural metadata:}
Catalog ID: \texttt{observer\_dependent\_truth}
Verified Tier: $O_2$
Crystal Address: 17279999 (approximate projection)
C-Score: 0.5505 - 0.5900
Promotions to $O_{\infty}$: 7


\end{document}
